% 2011 年安徽卷文科第 8 题：空间几何体的三视图，求表面积。
% 来源 PDF：普通高考/2011/2011安徽文.pdf，SHA-256 af8382bc2cd47a1d5e08bb622275f77ad100aa4f844b5cfd2fd5cd217de58ffc。
% 原页：第 1 页（A3 三栏）第 1 栏中部，书页码 645。
% 重绘依据：照原页——左上为正（主）视图，边长 4 的正方形，上方标 4、左侧标 4；
%   右上为侧（左）视图，下底 4、上底 2、高 4 的等腰梯形，与主视图上下边齐平；
%   左下为俯视图，边长为 4 的正方形，内部有两条水平线把高分成 1、2、1，
%   左侧以三个双向箭头分别标出 1、2、1；三视图下各标名称，原页不画顶点圆点、不标其他尺寸。
\begin{tikzpicture}[scale=0.55, >=stealth, semithick]
  % 正（主）视图：4 x 4 正方形
  \draw (0,0) -- (4,0) -- (4,4) -- (0,4) -- cycle;
  \draw[<->] (-0.05,4.55) -- (4.05,4.55);
  \node at (2,4.72) {$4$};
  \draw[<->] (-0.55,0) -- (-0.55,4);
  \node at (-0.78,2) {$4$};
  \node at (2,-0.62) {正（主）视图};
  % 侧（左）视图：下底 4、上底 2、高 4 的等腰梯形
  \draw (5.6,0) -- (9.6,0) -- (8.6,4) -- (6.6,4) -- cycle;
  \node at (7.6,-0.62) {侧（左）视图};
  % 俯视图：4 x 4 正方形，两条水平线分成 1、2、1
  \draw (0,-1.4) rectangle (4,-5.4);
  \draw (0,-2.4) -- (4,-2.4);
  \draw (0,-4.4) -- (4,-4.4);
  \draw[<->] (-0.55,-1.4) -- (-0.55,-2.4);
  \node at (-0.78,-1.9) {$1$};
  \draw[<->] (-0.55,-2.4) -- (-0.55,-4.4);
  \node at (-0.78,-3.4) {$2$};
  \draw[<->] (-0.55,-4.4) -- (-0.55,-5.4);
  \node at (-0.78,-4.9) {$1$};
  \node at (2,-6.0) {俯视图};
\end{tikzpicture}
